% Standalone problem document, generated from the adjacent README.md.
% Compile with XeLaTeX. Mathematical content is maintained in Markdown.
\documentclass[11pt,a4paper]{article}
\usepackage[fontsize=10.5pt]{scrextend}
\usepackage[margin=22mm,headheight=15pt,headsep=9mm,footskip=11mm]{geometry}
\usepackage{fontspec}
\setmainfont{texgyrepagella}[Extension=.otf,UprightFont=*-regular,BoldFont=*-bold,ItalicFont=*-italic,BoldItalicFont=*-bolditalic]
\setsansfont{texgyreheros}[Extension=.otf,UprightFont=*-regular,BoldFont=*-bold,ItalicFont=*-italic,BoldItalicFont=*-bolditalic]
\setmonofont{texgyrecursor}[Extension=.otf,UprightFont=*-regular,BoldFont=*-bold,ItalicFont=*-italic,BoldItalicFont=*-bolditalic,Scale=MatchLowercase]
\usepackage{amsmath,amssymb}
\usepackage{unicode-math}
\setmathfont{texgyrepagella-math.otf}
\usepackage{xcolor,microtype,enumitem,needspace,fancyhdr}
\usepackage{bookmark}
\definecolor{ink}{HTML}{17344B}
\definecolor{accent}{HTML}{197D80}
\definecolor{muted}{HTML}{596773}
\hypersetup{colorlinks=true,linkcolor=accent,urlcolor=accent,pdftitle={NR-03 - Full
nonnegative rank of the quadratic correlation
matrix},pdfauthor={Open Problems in Numerical Linear Algebra}}
\urlstyle{same}
\setlength{\parindent}{0pt}
\setlength{\parskip}{5pt plus 1pt minus 1pt}
\linespread{1.04}
\setlength{\emergencystretch}{3em}
\widowpenalty=10000
\clubpenalty=10000
\displaywidowpenalty=10000
\allowdisplaybreaks[1]
\setlist{nosep,leftmargin=1.5em}
\providecommand{\tightlist}{\setlength{\itemsep}{0pt}\setlength{\parskip}{0pt}}
\providecommand{\pandocbounded}[1]{#1}
\pagestyle{fancy}
\fancyhf{}
\fancyhead[L]{\sffamily\footnotesize\color{muted}OPEN PROBLEMS IN NUMERICAL LINEAR ALGEBRA}
\fancyhead[R]{\sffamily\bfseries\footnotesize\color{accent}NR-03}
\fancyfoot[L]{\parbox[b]{0.9\textwidth}{\sffamily\fontsize{8}{10}\selectfont\color{muted}Editorial ratings; a bounded search cannot certify that a problem remains unsolved.\\Literature check: 2026-09-10}}
\fancyfoot[R]{\sffamily\footnotesize\color{muted}\thepage}
\renewcommand{\headrulewidth}{0.3pt}
\renewcommand{\headrule}{\hbox to\headwidth{\color{accent}\leaders\hrule height \headrulewidth\hfill}}
\makeatletter
\renewcommand{\section}{\Needspace{5\baselineskip}\@startsection{section}{1}{0pt}{12pt plus 2pt minus 2pt}{4pt}{\sffamily\bfseries\large\color{ink}}}
\renewcommand{\subsection}{\Needspace{4\baselineskip}\@startsection{subsection}{2}{0pt}{9pt plus 1pt minus 1pt}{3pt}{\sffamily\bfseries\normalsize\color{ink}}}
\makeatother
\setcounter{secnumdepth}{0}
\begin{document}
{\sffamily\small\color{muted}Nonnegative and positive
factorizations\par}
\vspace{3pt}
{\raggedright\hyphenpenalty=10000\exhyphenpenalty=10000\sffamily\bfseries\fontsize{21}{24}\selectfont\color{ink}Full
nonnegative rank of the quadratic correlation matrix\par}
\vspace{7pt}
{\sffamily\small\textbf{Difficulty:} extreme\quad\textbf{Importance:} broadly
interesting\par}
{\sffamily\small\bfseries\color{accent}Status: Open\par}
\vspace{4pt}
{\color{accent}\hrule height 0.8pt}
\vspace{6pt}
\textbf{Rating rationale:} Extreme because full nonnegative rank is much
sharper than available exponential lower bounds for this family; broad
importance concerns communication complexity and limitations of linear
programming formulations.

\textbf{Area:} exact NMF and lower bounds for optimization formulations

\subsection{Context and notation}\label{context-and-notation}

All factorizations are over the real numbers. For
\(X\in\mathbb R_{\ge0}^{m\times n}\), define

\[
\operatorname{rank}_+(X)=\min\{r\ge0:X=WH,\quad
W\in\mathbb R_{\ge0}^{m\times r},\ H\in\mathbb R_{\ge0}^{r\times n}\}.
\]

\subsection{Problem statement}\label{problem-statement}

For each integer \(n\ge3\), let \(C_n\) be the \(2^n\times2^n\) matrix
indexed by \(a,b\in\{0,1\}^n\) and defined by

\[
C_n(a,b)=(1-a^{\mathsf T}b)^2.
\]

\subsubsection{Question}\label{question}

Is \(\operatorname{rank}_+(C_n)=2^n\) for every \(n\ge3\)? All entries,
including those for \(a^{\mathsf T}b>1\), are fixed by this formula.

The matrix has ordinary rank \(1+n(n+1)/2\) and is a submatrix of a
slack matrix of the correlation polytope formed from valid, possibly
redundant inequalities. The conjecture therefore asks for a large
separation between ordinary and nonnegative matrix rank with
consequences for linear programming representations.

\subsection{References}\label{references}

Vandaele, Gillis, Glineur, and Tuyttens,
\href{https://arxiv.org/html/1411.7245}{\emph{Heuristics for Exact
Nonnegative Matrix Factorization}}, §6.4, Conjecture 4. Gillis,
\href{https://orbi.umons.ac.be/bitstream/20.500.12907/42337/1/NMFbook_SIAM_reprint.pdf}{\emph{Nonnegative
Matrix Factorization}}, §3.7, p.~96, using the name \(U_n\) for this
fixed matrix.

\subsection{Status check --- 2026-09-10}\label{status-check-2026-09-10}

Rechecked \href{https://arxiv.org/html/2605.14058v2}{Baeckelant et
al.~v2, §6.7 and Appendix A.4} and
\href{https://arxiv.org/abs/2607.27014}{Sergeev's July 2026 paper}, and
searched for a full-rank resolution. The former retains the
prescribed-completion conjecture, including the unresolved n=3 case.
Sergeev treats a partial unique-disjointness matrix whose remaining
entries may vary; this does not determine the rank of the fixed matrix
here. No full resolution was located.
\end{document}
